\section{Block architecture}
\label{sec:arch}

This section is the paper's main artifact. It defines an RDMA endpoint in
nineteen blocks, targeting an AMD Alveo U55C accelerator card first and
expandable to an application-specific implementation, such that every block
names the measured behaviour it is required to reproduce, the evidence class of
every parameter it carries, the block of the golden C model it corresponds to,
and the control and status registers it publishes.

\subsection{Target}
\label{sec:arch:target}

The first target is one Alveo U55C card~\cite{alveo-u55c}.
Table~\ref{tab:u55c} lists the resources the architecture depends on. The choice is deliberate rather than
convenient: the card carries two 100\,GbE-capable cages and the integrated
Ethernet hard blocks to drive them, which is exactly the port count the
measured device has; it carries high-bandwidth memory large enough that
queue-pair context and memory translation need not be host resident at the
scale we care about; and its PCIe attachment is the same generation as the
measured host, so a comparison against measured host-side behaviour is not
confounded by the link.

\begin{table}[t]
  \centering
  \small
  \caption{Alveo U55C resources the architecture depends on.
    \TODO{confirm every row against the current data sheet before submission;
    the numbers below are the ones the design was sized against.}}
  \label{tab:u55c}
  \begin{tabular}{@{}ll@{}}
    \toprule
    Resource & Figure used for sizing \\
    \midrule
    Logic & about 1.30\,M lookup tables, 2.6\,M registers \\
    On-chip memory & about 70\,Mb block RAM, 270\,Mb ultra RAM \\
    Device memory & 16\,GB HBM2, about 460\,GB/s aggregate \\
    Network & 2 QSFP28 cages, 2 integrated 100\,GbE MAC cores \\
    Host & PCIe Gen3 $\times 16$, or Gen4 $\times 8$ \\
    \bottomrule
  \end{tabular}
\end{table}

Every dimension that a different target would change is a parameter, listed in
\S\ref{sec:arch:expand}: link rate, path maximum transmission unit, queue-pair
count, memory-region count, buffer sizes and port count. The application-specific
path replaces high-bandwidth memory with on-die static RAM and replaces the
soft interconnect adapters with hard macros, and keeps the same
register-transfer-level sources and the same register map.

\subsection{Two interconnects}
\label{sec:arch:interconnect}

The architecture separates moving data from observing the machine, and gives
each its own interconnect.

\textbf{The AXI4 data path} carries everything with a byte count: host-memory
reads for work-queue elements, context, translations and payload; host-memory
writes for placed payload and completion-queue entries; and high-bandwidth
memory traffic for context, translation and buffer backing. It is a
512-bit AXI4 fabric, which at the Ethernet core clock of 322.266\,MHz supplies
about 165\Gbps{} per direction, comfortably above the 100\Gbps{} line rate and
above the 197\Gbps{} internal budget the measurements found when loopback and
wire traffic share the machine. Read and write channels are independent, and
every initiator carries a transaction identifier that maps to one of the
transaction classes the model already accounts separately: user access region,
doorbell record, work-queue element, context, translation, payload read,
payload write, completion-queue entry, command and interrupt.

\textbf{The AXI4-Lite monitoring network on chip} carries nothing with a byte
count. It is a unidirectional ring of register slices, one stop per block, with
a single initiator at the host interface, 32-bit data and a 4\,KiB aperture per
stop. A ring rather than a crossbar because a crossbar with nineteen slaves and
one master is all cost and no benefit here: monitoring traffic is rare, entirely
tolerant of latency, and the ring's one register stage per stop is free timing
closure and a fixed, small resource cost that does not grow quadratically as
blocks are added. The ring is physically and logically separate from the data
path, so reading a counter cannot perturb the datapath timing that the counter
describes, which is the property that makes the resulting traces usable as
verification evidence.

\subsection{The register map}
\label{sec:arch:csr}

Every block occupies one 4\,KiB page of a single versioned map, addressed from
the control base address register. The first sixteen bytes of every page are
identical:

\begin{center}\small
\begin{tabular}{@{}lll@{}}
  \toprule
  Offset & Field & Meaning \\
  \midrule
  \texttt{0x000} & \csr{MAP\_VERSION}   & major and minor of the whole map \\
  \texttt{0x004} & \csr{BLOCK\_ID}      & B1 to B19 \\
  \texttt{0x008} & \csr{BLOCK\_VERSION} & this block's own version \\
  \texttt{0x00C} & \csr{CAPABILITIES}   & optional features present \\
  \bottomrule
\end{tabular}
\end{center}

\noindent
Software that knows \csr{MAP\_VERSION} knows the layout of every page, and a
block that gains a field bumps \csr{BLOCK\_VERSION} without invalidating
readers of the others. Three map-wide mechanisms follow.

\textbf{One clock.} A free-running 64-bit counter, clocked from the reference
the Ethernet subsystem already needs, ticks every 6.4\nS{} at 156.25\,MHz. It is
broadcast to every block, and every timestamped event in the design is stamped
from it. The measured device offers a hardware timestamp on received packets
and nothing else; here every block's events share one time base.

\textbf{Coherent snapshots.} Writing \csr{SNAPSHOT} at the map level asserts a
broadcast strobe that latches every block's counters and occupancies into a
shadow bank together with the current timestamp. The campaign spent real effort
working around the fact that counters on the measured device are read
non-atomically and are cached by the driver for 10\mS{}, which is why rates
below that window are meaningless there. Here a read set is coherent by
construction.

\textbf{Counter width and a compatibility mode.} All counters are 64 bits, which
removes the wrap that makes \ctr{rx\_write\_requests} unusable at high message
rates on the real device. Because fidelity to the real device is sometimes the
point, a \csr{WRAP32\_COMPAT} bit makes the mirrored counters wrap at 32 bits
exactly as the original does, so a tool written against the real adapter can be
tested against this one.

The map has two halves. The \emph{mirror} half carries counters named exactly as
the real device names them, with the real semantics preserved including the ones
that are wrong: \ctr{np\_ecn\_marked\_roce\_packets} is present and stays zero
(\anom{07}), \ctr{packet\_seq\_err} counts loss bursts and not lost packets
(\anom{19}), \ctr{out\_of\_buffer} exists and never moves under the conditions
the campaign covered, and a \csr{FW\_VARIANT} field selects whether
\ctr{local\_ack\_timeout\_err} counts or stays zero (\anom{09}). The
\emph{extension} half carries what the real device does not expose and the FPGA
can: the instantaneous and high-water occupancy of every internal queue, the
credit state of every arbiter, per-stage service counts, and the reaction
point's rate and $\alpha$ per queue pair. \S\ref{sec:verify:fpga} returns to
what that buys.

\begin{figure*}[t]
  \centering
  %% Top-level block architecture.  Pure TikZ, no external assets.
%% Coordinates are in units of 0.9 cm; the picture is sized to the text width
%% of a two-column page and does not depend on the document class.
\begin{tikzpicture}[
  x=0.9cm, y=0.9cm,
  font=\scriptsize,
  blk/.style   = {draw, rounded corners=2pt, align=center, fill=black!4,
                  text width=18mm, inner sep=2pt, minimum height=7.5mm},
  tall/.style  = {draw, rounded corners=2pt, align=center, fill=black!4,
                  text width=15mm, inner sep=2pt, minimum height=48mm},
  ext/.style   = {align=center, font=\scriptsize},
  bus/.style   = {draw, fill=black!12, rounded corners=1pt},
  flow/.style  = {-{Latex[length=1.7mm,width=1.5mm]}, semithick},
  duplex/.style= {{Latex[length=1.7mm,width=1.5mm]}-%
                  {Latex[length=1.7mm,width=1.5mm]}, semithick},
  axi/.style   = {-{Latex[length=1.4mm,width=1.3mm]}, thin, draw=black!65},
  mon/.style   = {dashed, draw=black!60, thin},
]

%% ==== monitoring ring, drawn first so blocks sit on top of it =============
\draw[mon, rounded corners=10pt] (0.3,-1.0) rectangle (18.2,6.5);

%% ==== host and wire, outside the ring =====================================
\node[ext] (host) at (1.6,7.2) {Host\\PCIe Gen3 $\times$16 / Gen4 $\times$8};
\node[ext] (wire) at (16.7,7.2) {Wire\\2 $\times$ 100\,GbE};

%% ==== left and right edge blocks ==========================================
\node[tall] (b1) at (1.6,3.1)
  {\textbf{B1}\\Host interface\\and BAR map\\[3pt]UAR doorbell pages,
   CSR BAR, MSI-X};
\node[tall] (b9) at (16.7,3.1)
  {\textbf{B9}\\MAC / PHY\\[3pt]100\,GbE cores, global pause emission};

%% ==== lane 1: transmit datapath ===========================================
\node[blk] (b2)  at (4.6,5.4)  {\textbf{B2} Submission\\SQ/RQ, doorbell};
\node[blk] (b3)  at (7.2,5.4)  {\textbf{B3} WQE fetch\\DMA, prefetch};
\node[blk] (b6)  at (9.8,5.4)  {\textbf{B6} Scheduler\\RR, window, pacer};
\node[blk] (b7)  at (12.4,5.4) {\textbf{B7} Packetizer\\headers, ECT(0),
  ICRC};

%% ==== lane 2: context, translation, gather, loopback ======================
\node[blk] (b4)  at (4.6,4.0)  {\textbf{B4} QPC\\cache $+$ HBM};
\node[blk] (b5)  at (7.2,4.0)  {\textbf{B5} MR / MTT\\MPT, page walks};
\node[blk] (b8)  at (9.8,4.0)  {\textbf{B8} TX DMA\\gather, MRRS};
\node[blk] (b17) at (12.4,4.0) {\textbf{B17} Loopback\\internal budget};

%% ==== the AXI4 data path bus ==============================================
\draw[bus] (3.4,3.04) rectangle (15.0,3.36);
\node[font=\scriptsize] at (9.2,3.2) {AXI4 data path: host DMA and HBM};

%% ==== lane 3: management and congestion control ===========================
\node[blk] (b18) at (5.9,2.4)  {\textbf{B18} Management\\GID, MTU, QP
  transitions};
\node[blk] (b15) at (11.1,2.4) {\textbf{B15} Reaction pt\\DCQCN rate,
  $\alpha$};
\node[blk] (b14) at (13.7,2.4) {\textbf{B14} Notif.\ point\\CNP generation};

%% ==== lane 4: receive datapath ============================================
\node[blk] (b16) at (4.6,1.0)  {\textbf{B16} Completion\\CQE, EQ, MSI-X};
\node[blk] (b13) at (7.2,1.0)  {\textbf{B13} Requester\\ACK, go-back-N};
\node[blk] (b12) at (9.8,1.0)  {\textbf{B12} Responder\\PSN, ACK, placement};
\node[blk] (b11) at (12.4,1.0) {\textbf{B11} Parser\\classify, ICRC};
\node[blk] (b10) at (15.0,1.0) {\textbf{B10} Ingress meter\\PHY discard};

%% ==== transmit flow =======================================================
\draw[flow] (host) -- (b1.north);
\draw[flow] (b1.east|-b2.west) -- (b2.west);
\draw[flow] (b2) -- (b3);
\draw[flow] (b3) -- (b6);
\draw[flow] (b6) -- (b7);
\draw[flow] (b7.east) -| ([xshift=-4mm]b9.north);
\draw[duplex] ([xshift=4mm]b9.north) -- (wire);

%% ==== receive flow ========================================================
\draw[flow] (b9.south) |- (b10.east);
\draw[flow] (b10) -- (b11);
\draw[flow] (b11) -- (b12);
\draw[flow] (b12) -- (b13);
\draw[flow] (b13) -- (b16);
\draw[flow] (b16.west) -- (b1.east|-b16.west);

%% ==== loopback and the shared ingress budget ==============================
\draw[flow] (b7.south) -- (b17.north);
\draw[flow] (b17.east) -| ([xshift=8mm]b10.north);

%% ==== congestion control ==================================================
\draw[flow] (b10.north) -- (b14.south);
\draw[flow] (b11.north west) -- (b15.south east);
\draw[flow] (b15.west) -- ++(-0.55,0)
  node[left, font=\scriptsize] {rate to B6};

%% ==== AXI4 data path attachments ==========================================
\draw[axi] (b1.east) -- (3.4,3.2);
\draw[axi] (b4.south)  -- (4.6,3.36);
\draw[axi] (b5.south)  -- (7.2,3.36);
\draw[axi] (b8.south)  -- (9.8,3.36);
\draw[axi] (b17.south) -- (12.4,3.36);
\draw[axi] (b16.north) -- (4.6,3.04);
\draw[axi] (b12.north) -- (9.8,3.04);
\draw[axi] (b3.east) -- (8.5,5.4) -- (8.5,3.36);

%% ==== the monitoring ring endpoints =======================================
\draw[mon] (b1.west)  -- (0.3,3.1);
\draw[mon] (b9.east)  -- (18.2,3.1);
\draw[mon] (b2.north) -- (4.6,6.5);
\draw[mon] (b6.north) -- (9.8,6.5);
\draw[mon] (b7.north) -- (12.4,6.5);
\draw[mon] (b16.south) -- (4.6,-1.0);
\draw[mon] (b12.south) -- (9.8,-1.0);
\draw[mon] (b10.south) -- (15.0,-1.0);
\node[blk, fill=white, text width=24mm] (b19) at (2.4,-1.0)
  {\textbf{B19} CSR, counters, timestamp};
\node[font=\scriptsize, anchor=west] at (10.6,-1.4)
  {AXI4-Lite monitoring ring, one stop per block};

\end{tikzpicture}

  \caption{Top-level block architecture. Solid arrows are the AXI4 data path
    and the packet path; the dashed ring is the AXI4-Lite monitoring network on
    chip, which reaches every block and is the only path by which state leaves
    the design for observation. Blocks are numbered as in
    \S\ref{sec:arch:blocks} and in the block column of
    Table~\ref{tab:anomaly}.}
  \Description{Top-level block diagram of the RDMA endpoint: a transmit lane, a context and translation lane, an AXI4 data path bus, a congestion-control lane and a receive lane, between a host interface on the left and a media access block on the right, all enclosed by a dashed AXI4-Lite monitoring ring.}
  \label{fig:arch}
\end{figure*}

\subsection{The blocks}
\label{sec:arch:blocks}

Each block below states, in order: what it does, the state it owns, its
interfaces, the measured behaviour it is required to reproduce with the finding
and anomaly row that pin it, the corresponding block of the golden C model, and
the registers it publishes on the monitoring ring. Register names written as
\ctr{this} are mirrors of names the real device uses; names written as
\csr{THIS} belong to this design.

\subsubsection{B1: Host interface and BAR map}

\textbf{Function.} Terminate the PCIe link and split its traffic in two. Base
address register 0 is the control region: it decodes into the AXI4-Lite
monitoring ring and is the only initiator on it. Base address registers 2 and 4
expose user access region pages, one page per doorbell context, mapped
write-combining into user processes; a store to such a page is a doorbell, and
a store of a full work-queue element to such a page is the inline publication
path that lets the device skip the fetch entirely. The block also owns the
MSI-X table and the pending-bit array, and converts posted writes, non-posted
reads and completions into AXI4 transactions tagged with the class they belong
to.

\textbf{State.} Base-address decode; the user-access-region page ownership map;
write-combining coalescing buffers; the MSI-X table and pending bits;
outstanding read tags and the completion reservation for them.

\textbf{Interfaces.} PCIe Gen3 $\times 16$ or Gen4 $\times 8$ toward the host;
AXI4 master and slave toward the data path; AXI4-Lite master onto the
monitoring ring; a doorbell stream to B2.

\textbf{Must reproduce.} That PCIe is not the bottleneck at 100\,GbE. The
campaign's stall counters were exactly zero at every one of more than 700
points, including 195\Gbps{} of host-loopback traffic at 4\,MiB messages across
64 queue pairs, and including the 197\Gbps{} internal-budget point. The block
must therefore expose \ctr{outbound\_pci\_stalled\_rd} and
\ctr{outbound\_pci\_stalled\_wr} and they must be capable of moving, so that a
later study at a higher rate can show when the link does bind; a model in which
they are hard-wired to zero would be unfalsifiable.

\textbf{Evidence.} The base-address layout, the user-access-region semantics and
the doorbell record are \evdoc{}. Write-combining behaviour and the inline
publication path are \evdrv{} from the provider source. Maximum payload size,
maximum read request size, read-completion-boundary policy and credit pools are
\evdoc{}. The service latencies of each transaction class are \evcal{}.

\textbf{Golden model.} The transaction-level PCIe fabric with its per-class
inventory (user access region, BlueFlame, doorbell record, work-queue element,
context, translation, payload read, payload write, completion-queue entry,
command, interrupt) and the work-queue PCIe binding.

{\raggedright
\textbf{Registers.} \csr{PCI\_LINK\_GEN}, \csr{PCI\_LINK\_WIDTH}, \csr{MPS},
\csr{MRRS}, \ctr{outbound\_pci\_stalled\_rd},
\ctr{outbound\_pci\_stalled\_wr}, \csr{UAR\_WRITES},
\csr{INLINE\_PUBLICATIONS}, \csr{DB\_RECORD\_READS}, and, in the extension
half, \csr{READ\_TAG\_OCCUPANCY} and \csr{WC\_BUFFER\_FLUSHES}.\par}

\subsubsection{B2: Submission engine}

\textbf{Function.} Own the send and receive queue rings. On a doorbell, read the
doorbell record, compute the accepted prefix of newly posted work-queue
elements, and hand their indices to B3 as one batch. Handle inline data, which
arrives with the doorbell rather than through a fetch. Retire slots when B16
says a completion has made them reclaimable.

\textbf{State.} Per queue pair: ring base address and logarithmic size, the
producer index last read from the doorbell record, the device's own consumer
index, the accepted-prefix pointer, and an inline staging buffer.

\textbf{Interfaces.} Doorbell stream from B1; AXI4 read for doorbell records;
work-queue-element index batches to B3; reclaim notifications from B16.

\textbf{Must reproduce.} Accepted-prefix semantics: a single doorbell publishes
an arbitrary prefix of posted work, so the number of memory-mapped writes per
work request falls with the posting batch. The measured consequence is that
posting batch is nearly free: batches of 1 and 64 differ by only $1.16\times$,
because both arms pipeline at the same queue depth. The default send-queue
depth of 1024 must be a parameter and not an assumption, because the campaign's
single most expensive analysis error was reading a tool flag that sets work
requests per posting call as if it set queue depth.

\textbf{Evidence.} Ring geometry, doorbell records and inline semantics are
\evdoc{} and \evdrv{}. The depth default is \evdrv{}. The $1.16\times$ batching
figure is a measured property of the pipeline, and belongs to B6.

\textbf{Golden model.} The submission model and the work-queue core's
accepted-prefix posting with explicit doorbell batches.

{\raggedright
\textbf{Registers.} \csr{SQ\_DEPTH}, \csr{RQ\_DEPTH}, \csr{DOORBELLS\_SEEN},
\csr{WQES\_ACCEPTED}, \csr{INLINE\_BYTES}; extension half
\csr{ACCEPTED\_PREFIX\_HIST}, \csr{SQ\_OCCUPANCY},
\csr{SQ\_OCCUPANCY\_HIGH\_WATER}.\par}

\subsubsection{B3: Work-queue-element fetch}

\textbf{Function.} Read work-queue-element bytes from host memory by direct
memory access, prefetching across a doorbell batch so that the fetch of one
element overlaps the processing of the previous one. Bypassed entirely when the
element arrived inline.

\textbf{State.} Outstanding fetch tags; the prefetch window; a per-queue-pair
fetch cursor; a small reorder buffer, because read completions may return out
of order.

\textbf{Interfaces.} AXI4 read to host memory; decoded work-queue-element
records to B4 and B6.

\textbf{Must reproduce.} One component of the fixed-offset initiation cost
$T_{\text{eff}} = 4.48$\uS{}. The campaign can identify the sum of the
initiation stages and not their split, so the architecture assigns a
parameterised service to each of the five stages (publication, observation,
fetch, context readiness, admission) whose sum is constrained and whose split
is \evcal{}. The block must also not become the depth limiter: with 1024
elements outstanding the machine runs $5.9\times$ faster at 8\,KiB than with
one, so prefetch depth is sized from the outstanding-work window and not the
other way round.

\textbf{Evidence.} The element layout is \evdoc{}; the fetch batching is
\evdrv{}; the stage split of $T_{\text{eff}}$ is \evcal{} and explicitly an
assumption until a work-queue publication campaign lands.

\textbf{Golden model.} The serialized fetch stage of the queue core, backed by
read and completion transactions on the PCIe fabric.

{\raggedright
\textbf{Registers.} \csr{WQE\_FETCH\_BYTES}, \csr{WQE\_FETCH\_COUNT},
\csr{PREFETCH\_DEPTH}; extension half \csr{WQE\_FETCH\_OUTSTANDING},
\csr{WQE\_FETCH\_LATENCY\_HIST}.\par}

\subsubsection{B4: Queue-pair context}

\textbf{Function.} Hold the per-queue-pair state that both the transmit and
receive paths read, in an on-chip cache backed by high-bandwidth memory, and
serve the lookup stage of both paths.

\textbf{State.} Per queue pair: transport type; state in the RESET, INIT, RTR,
RTS, SQD, SQE, ERR machine; send packet sequence number and expected receive
packet sequence number; last acknowledged sequence number; path maximum
transmission unit; destination queue-pair number, global identifier, MAC and
VLAN; traffic class and differentiated-services value; local acknowledgement
timeout and retry counts; the receiver-not-ready timer; the outstanding-work
window; and the congestion-control state, namely the reaction point's $\alpha$,
current rate, target rate, stage, timer and byte counter, and the notification
point's last-CNP timestamp.

\textbf{Interfaces.} AXI4 to high-bandwidth memory for the backing store;
lookup ports to B6 on the transmit side and B11 on the receive side; update
ports from B12, B13 and B15; configuration writes from B18.

\textbf{Must reproduce.} That congestion-control state is persistent per queue
pair and survives across work requests. This is what makes \anom{18} possible:
a rate cut that takes 3 to 39\mS{} to arrive and $447 \pm 10$\mS{} to recover
cannot be modelled by state that resets per message. Path maximum transmission
unit lives here and drives \anom{13}, the 5.6\,\% goodput difference between
1024 and 4096. The state machine's transitions and their legality are
\evdoc{}.

\textbf{Evidence.} Every field is \evdoc{} from the reference manual or the
specification. The cache geometry is \evcal{} with a measured lower bound: 480
active queue pairs against 12\,000 memory regions produced no throughput effect
at all (\anom{15}), so the real device's working set at that scale fits, and
the associativity and replacement policy are not identifiable from outside.

\textbf{Golden model.} The context lookup stage, plus the transport and
congestion-control field groups of the hardware profile.

{\raggedright
\textbf{Registers.} \csr{QP\_COUNT\_ACTIVE}, \csr{QPC\_HITS},
\csr{QPC\_MISSES}, \csr{QP\_SELECT} with a per-queue-pair window exposing
\csr{QP\_STATE}, \csr{QP\_PSN\_SEND}, \csr{QP\_PSN\_EXPECTED} and
\csr{QP\_MTU}; extension half \csr{QPC\_MISS\_SERVICE\_NS},
\csr{QP\_STATE\_HIST}, \csr{QP\_RATE}, \csr{QP\_ALPHA}.\par}

\subsubsection{B5: Memory registration and translation}

\textbf{Function.} Hold the memory protection table (region base, length, keys,
protection domain, access rights) and the memory translation table
(virtual-to-physical mapping), cache both on chip, validate keys and access
rights on every reference, and walk host memory on a miss.

\textbf{State.} Cached protection and translation entries; the page-walk queue;
key and access-right check results.

\textbf{Interfaces.} Lookup ports from B8 (transmit gather) and B12 (receive
placement); AXI4 read for walks; configuration writes from B18.

\textbf{Must reproduce.} \anom{15}, and it is the one row of the table that is
reproduced \emph{by absence}. Twelve thousand memory regions cost $+1.2\,\%$
against a two-region control, and 1024 regions of 64\,KiB were identical to the
control to four significant figures. The architecture therefore sizes the cache
above the measured working set and the block adds no cost in that regime; what
it must \emph{not} do is silently hide a smaller cache, which is why the miss
counters and the service time a miss cost are published separately. The real
device's behaviour beyond the measured working set is unknown and is declared
so.

\textbf{Evidence.} Table layouts, keys and access rights are \evdoc{}; the
cache size is \evcal{} with a lower bound; behaviour past the bound is
\evdecl{}.

\textbf{Golden model.} Not modelled in the current slice, and documented as an
absence rather than assumed to be free; the profile's three-tier context
hierarchy (on-die static RAM, optional device memory, host pinned memory) is
where it lands.

{\raggedright
\textbf{Registers.} \csr{MR\_COUNT}, \csr{MPT\_HITS}, \csr{MPT\_MISSES},
\csr{MTT\_HITS}, \csr{MTT\_MISSES}, \csr{PAGE\_WALKS},
\csr{KEY\_VIOLATIONS}; extension half \csr{XLAT\_CACHE\_OCCUPANCY},
\csr{XLAT\_MISS\_SERVICE\_NS}.\par}

\subsubsection{B6: Scheduler and transmit arbiter}

\textbf{Function.} Choose which queue pair transmits next. Round robin over
queue pairs that are ready, subject to three gates: the per-queue-pair
outstanding-work window, which is released by acknowledgements; the
per-queue-pair and per-device packet-rate and bit-rate limiters; and the
per-queue-pair rate token supplied by the reaction point B15.

\textbf{State.} The ready list and its round-robin pointer; per queue pair the
outstanding bytes and packets and the token bucket; the device-wide token
bucket.

\textbf{Interfaces.} Work-queue-element records from B3; context lookup to B4;
rate from B15; window release from B13; packet-issue grants to B7.

\textbf{Must reproduce.} Four measurements pin this block.
(i)~Queue-depth dominance: 13.2\Gbps{} at depth one against 77.2\Gbps{} at
depth 1024, both at 8\,KiB, a factor of 5.9, which requires a genuine
outstanding-work window rather than one element at a time.
(ii)~\anom{14}, the host-bound message rate: 3.87 million packets per second on
one queue pair at 1\,KiB, and 16.7 million messages per second per sender at
512\,B, which the packet-rate limiters carry.
(iii)~Multi-queue-pair recovery of the ceiling: two or more queue pairs reach
97.7\Gbps{} where one saturating queue pair settles into the loss equilibrium.
(iv)~Skew costs completion time and not bandwidth: with seven queue pairs
carrying different message sizes, the per-queue-pair \emph{message} rate was
identical to three decimal places and the per-queue-pair \emph{byte} rate was
exactly proportional to message size (a spread of 2.000 for a 2:1 size ratio),
with an aggregate of 96.5\Gbps{}. That is precisely the signature of round
robin over message completions, and it is the strongest single piece of
evidence for the arbitration policy in this block.

\textbf{Evidence.} The arbitration policy is \evcal{}, inferred from (iv). The
rate-limiter values are \evcal{}. The existence of a per-queue-pair rate gate
fed by congestion control is \evdoc{} from the algorithm.

\textbf{Golden model.} The scheduler admission stage, the outstanding-work
window and the transmit pacer.

{\raggedright
\textbf{Registers.} \csr{TX\_PPS\_LIMIT}, \csr{TX\_BPS\_LIMIT},
\csr{SCHED\_GRANTS}, \csr{QP\_READY\_COUNT}; extension half
\csr{RR\_POINTER}, \csr{QP\_OUTSTANDING\_BYTES}, \csr{QP\_OUTSTANDING\_PKTS},
\csr{TX\_RATE\_TOKENS}.\par}

\subsubsection{B7: Packetizer and header builder}

\textbf{Function.} Turn a granted message into packets. Segment at the path
maximum transmission unit; allocate packet sequence numbers; build the base
transport header and whichever extended header the operation needs (RETH for
remote address and length, AETH for acknowledgement syndrome and credits, DETH
for the source queue pair of a datagram); wrap the result in the RoCEv2
encapsulation, choosing the UDP source port so that it is stable per queue pair
and spread across queue pairs; write the differentiated-services field through
unchanged and force the two ECN bits to ECT(0); and compute the invariant
cyclic redundancy check.

\textbf{State.} Per queue pair a message cursor and a sequence-number
allocator; header templates cached from the context; the check accumulator.

\textbf{Interfaces.} Grants from B6; context read from B4; payload beats from
B8; packets to B9 for the wire or B17 for loopback.

\textbf{Must reproduce.} \anom{13}, the 5.6\,\% maximum-transmission-unit tax,
which emerges from per-packet header overhead and needs no special rule.
\anom{06}, the ECT(0) stamp, which does not emerge from anything and is
therefore \emph{injected}: the block overwrites the two ECN bits
unconditionally on every RoCEv2 transmit while passing the
differentiated-services field through, exactly as the on-wire probe found. The
number of packets a message becomes here bounds the go-back-N replay of
\anom{11} from above (256 at 1\,MiB), while the measured amplification of
about 16 says the unit actually replayed is the requester's outstanding
window of about 64\,KiB; the incast tax is therefore a downstream consequence
of segmentation, the window and go-back-N, and not a separate rule. Finally
\anom{02}, the sequence-error storm under two-entry gathers, is injected here
as a per-packet drop rule keyed on gather-entry count and size, because the
mechanism is not public and no differential experiment we could run localized
it further than ``the transmit path with two-entry gathers is causal''.

\textbf{Evidence.} Header layouts, the check and the encapsulation are
\evdoc{}. The ECN overwrite is measured, and \evcal{} as to why. The
two-entry-gather rule is \evdecl{} and labelled injected in the table.

\textbf{Golden model.} The packetizer.

{\raggedright
\textbf{Registers.} \ctr{tx\_packets\_phy}, \csr{TX\_PAYLOAD\_BYTES},
\csr{TX\_HEADER\_BYTES}, \csr{PSN\_ALLOCATED}, \csr{ECT0\_STAMPED},
\csr{ICRC\_GENERATED}, \csr{MTU\_IN\_USE}; extension half
\csr{PACKETS\_PER\_MESSAGE\_HIST}.\par}

\subsubsection{B8: Transmit direct memory access}

\textbf{Function.} Fetch payload. Walk the scatter-gather list, translate each
entry through B5, segment reads at the maximum read request size, manage read
credits and tags, and reorder returning completions into the packet order B7
needs. Bypassed for inline data.

\textbf{State.} Per-packet gather cursor; outstanding read tags; the completion
reorder buffer; credit counters.

\textbf{Interfaces.} AXI4 read to host memory and to high-bandwidth memory for
device-resident buffers; translation lookups to B5; payload beats to B7.

\textbf{Must reproduce.} That gather entries are free. Two-entry gathers cost
nothing in throughput (the storm of \anom{02} is a loss effect, with goodput
within 3\,\%), and no gather configuration the campaign tried moved a
throughput number. That the reorder buffer never becomes the limit at
100\,GbE, consistent with the stall counters staying at zero.

\textbf{Evidence.} Maximum payload size, maximum read request size, the
read-completion-boundary policy and the credit pools are \evdoc{}. Reorder
buffer depth is \evcal{}.

\textbf{Golden model.} The payload-read transaction class of the PCIe fabric.

{\raggedright
\textbf{Registers.} \csr{TX\_DMA\_READ\_BYTES}, \csr{TX\_DMA\_READS},
\ctr{outbound\_pci\_stalled\_rd} (mirrored from B1); extension half
\csr{TX\_DMA\_OUTSTANDING}, \csr{REORDER\_OCCUPANCY}, \csr{SGE\_COUNT\_HIST}.\par}

\subsubsection{B9: Media access control and physical layer}

\textbf{Function.} Drive the two 100\,GbE ports through the device's integrated
Ethernet cores, and emit 802.3x global pause frames when B10 asks for them.
Priority flow control is implemented but disabled by default, to match the
deployment that was measured.

\textbf{State.} Link state per port; the pause state machine; the physical-layer
counter set.

\textbf{Interfaces.} Streaming interfaces to B7 and from B10; the optical cages;
a pause request line from B10.

\textbf{Must reproduce.} The endpoint half of \anom{08}. An overloaded receiver
\emph{does} emit global pause, on the order of 760 frames per twenty-second
run, and no endpoint in the entire campaign ever received a pause frame of any
kind: \ctr{rx\_pause\_ctrl\_phy} was exactly zero at every node in every phase.
The design must therefore never assume that backpressure propagates, and must
publish both directions of the pause counters so that a deployment where it
does propagate is distinguishable from one where it does not.

\textbf{Evidence.} \evdoc{} throughout (IEEE 802.3 and its pause annex).

\textbf{Golden model.} The flow-control field group of the hardware profile:
priority flow control disabled, global pause transmission enabled, pause
propagation false.

{\raggedright
\textbf{Registers.} \ctr{rx\_packets\_phy}, \ctr{tx\_packets\_phy},
\ctr{rx\_discards\_phy}, \ctr{tx\_discards\_phy}, \ctr{tx\_pause\_ctrl\_phy},
\ctr{rx\_pause\_ctrl\_phy}, \ctr{tx\_global\_pause},
\ctr{tx\_global\_pause\_duration}, \ctr{tx\_pause\_storm\_warning\_events},
\ctr{rx\_prio0\_bytes}, \ctr{tx\_prio0\_bytes}, \csr{LINK\_SPEED},
\csr{LINK\_UP}.\par}

\subsubsection{B10: Receive ingress meter}

\textbf{Function.} A finite buffer at the receive port, drained at a service
rate, into which every arriving packet is written before it is parsed. Overflow
discards at the physical layer, with no transport signal of any kind: the
sender learns about it only through the retransmission it eventually causes.
The block also raises the pause request that B9 turns into a frame, and the
congestion signal that B14 turns into a notification packet.

\textbf{State.} Occupancy and its high-water mark; the drain-rate token bucket;
the state of the stall process described below.

\textbf{Interfaces.} From B9 (wire) and B17 (loopback); to B11; pause request to
B9; congestion indication to B14.

\textbf{Must reproduce.} More of the campaign than any other block.
\anom{03}: a saturated deep pipeline settles at 78 to 92\Gbps{} of goodput with
the counter signature \ctr{rx\_discards\_phy} and \ctr{out\_of\_sequence} at
the receiver and \ctr{packet\_seq\_err} and \ctr{roce\_adp\_retrans} at the
sender.
\anom{04}: above a size-dependent inter-burst gap (at least 4\uS{} at 8\,KiB,
between 4 and 100\uS{} at 64\,KiB) the discards go to zero and the duty-cycle
model becomes exact to $\pm 0.1\,\%$; a 4\uS{} gap at 8\,KiB raises wall
goodput by 13.8\,\%.
\anom{10}: full duplex at 91.8\Gbps{} per direction is counter-clean while one
direction at 93.4\Gbps{} is not, so the threshold sits between them.
\anom{17}: a lone flow above about 94\Gbps{} loses $0.1798\,\% \pm 0.0012$ of
its packets here, in bursts of about 73 packets lasting about 94\uS{}, path
independent over 32 fresh five-tuples, with the event count falling
$12.5\times$ under at least 10\Gbps{} of reverse traffic while the burst length
does not move, and with \ctr{out\_of\_buffer} at zero throughout, so it is not
receive-queue exhaustion.

\textbf{Open clause.} The discard ledger does not close, and the paper says so
rather than hiding it. At the \anom{03} operating point the campaign counted
only a few thousand physical-layer discards over a thirty-second window, far
too few to explain a twenty-percent goodput deficit by loss alone. The meter
reproduces the equilibrium at a fitted drain rate; it does not explain it.
\anom{17} is the narrowed form of the missing mechanism, and a second anchor
makes the gap explicit: under fan-in the same receiver absorbed 99.39\Gbps{}
while discarding almost nothing, and one drain rate cannot satisfy both
anchors. This is the block with the most \evcal{} content in the architecture,
and the one whose first hardware build is most worth instrumenting, because the
monitoring ring can publish the occupancy trace that the real device cannot.

\textbf{Evidence.} The existence of a port ingress buffer is \evdoc{}; its
depth, its drain rate and the stall process are \evcal{}, and the stall process
is currently unmodelled.

\textbf{Golden model.} The ingress meter.

{\raggedright
\textbf{Registers.} \ctr{rx\_discards\_phy} (mirror), \ctr{out\_of\_buffer}
(mirror, inert), \csr{INGRESS\_DRAIN\_BPS}, \csr{INGRESS\_DEPTH\_BYTES},
\csr{STALL\_EVENTS}; extension half \csr{INGRESS\_OCCUPANCY},
\csr{INGRESS\_HIGH\_WATER}, \csr{STALL\_BURST\_LEN\_HIST},
\csr{PAUSE\_REQUESTS}.\par}

\subsubsection{B11: Receive parser and classifier}

\textbf{Function.} Decode Ethernet, IP, UDP and the base transport header;
verify the invariant cyclic redundancy check; look the destination queue pair
up in B4; and route the packet to the block that handles its class. Congestion
notification packets are a class of their own: they are a base transport header
with a distinguished opcode and no payload, and they go to B15, not to the
responder.

\textbf{State.} Pipeline registers; the classification table.

\textbf{Interfaces.} From B10; context lookup to B4; classified packets to B12
(requests and data), B13 (acknowledgements and negative acknowledgements), B15
(notifications) and B18 (management).

\textbf{Must reproduce.} Correct separation of the notification class,
without which the reaction point cannot be driven at all; the
congestion-experienced bit surfaced to B14 as a distinct input from B10's own
congestion signal, because the measured deployment exercises the second and
never the first; and exact counting of check failures and of packets whose
queue-pair number matches nothing.

\textbf{Evidence.} \evdoc{} throughout.

\textbf{Golden model.} The front half of the receive processor.

{\raggedright
\textbf{Registers.} \csr{RX\_PACKETS}, \csr{RX\_BYTES}, \csr{ICRC\_ERRORS},
\csr{UNMATCHED\_QPN}, \csr{CNP\_RECEIVED}, \csr{BAD\_HEADER}; extension half
\csr{CLASS\_HIST}, \csr{CE\_MARKED\_RECEIVED}.\par}

\subsubsection{B12: Responder}

\textbf{Function.} The receive side of the transport, and the block that
decides whether an arriving packet is placed or thrown away. Check the packet
sequence number against the expected one; on a match, place the payload by direct memory access
for a write, consume a receive-queue entry for a send, or queue a read
response; on a gap, emit a negative acknowledgement naming the first missing
number and discard until the requester restarts. Generate acknowledgements,
with coalescing. Handle receiver-not-ready when no receive buffer is posted.

\textbf{State.} Expected sequence number per queue pair; acknowledgement
coalescing state; the read-response queue; the receiver-not-ready timer; the
receive-queue consumer index.

\textbf{Interfaces.} Classified packets from B11; context read and update to
B4; AXI4 write to host memory for placement; translation lookup to B5;
acknowledgement, negative-acknowledgement and read-response packets to B7;
completion notifications to B16.

\textbf{Must reproduce.} \ctr{out\_of\_sequence} moving under exactly the
conditions measured (about 93\,000 per thirty seconds under two-to-one incast,
zero for paced traffic). \ctr{rx\_write\_requests} incrementing exactly once
per completed RDMA write message, which the campaign verified against sender
completion accounting to 0.01\,\% and used as its ground-truth goodput counter,
and wrapping at 32 bits when the compatibility bit is set. And the negative
result: receiver-not-ready never fired and \ctr{out\_of\_buffer} never moved in
the entire campaign, because receive queues were posted deep and RDMA writes
consume none, so overflow appears one layer down as an ingress discard rather
than as a transport event.

\textbf{Evidence.} \evdoc{} (the specification fixes the sequence-number check,
the syndromes and the recovery rule). Acknowledgement coalescing is \evcal{}.

\textbf{Golden model.} The responder half of the receive processor.

{\raggedright
\textbf{Registers.} \ctr{rx\_write\_requests}, \csr{RX\_READ\_REQUESTS},
\ctr{out\_of\_sequence}, \ctr{implied\_nak\_seq\_err},
\ctr{rnr\_nak\_retry\_err}, \ctr{out\_of\_buffer}, \csr{ACKS\_SENT},
\csr{NAKS\_SENT}, \csr{READ\_RESPONSES\_SENT}.\par}

\subsubsection{B13: Requester completion and retransmission}

\textbf{Function.} The send-side transport. Process acknowledgements and
negative acknowledgements, advance the unacknowledged window and release it to
B6, replay from the first unacknowledged sequence number on a negative
acknowledgement or a local acknowledgement timeout, and count retries.
Go-back-N is the default; a limited selective-repeat window is a parameter and
is off in the ConnectX-5 profile.

\textbf{State.} Per queue pair the unacknowledged window and its replay cursor;
retry counters; the retransmission timer.

\textbf{Interfaces.} From B11; context update to B4; replay requests to B6 and
B7; window release to B6; completions to B16.

\textbf{Must reproduce.} The amplification that turns a small loss rate into a
large goodput deficit. \anom{11}: 1.65\,\% packet loss becomes a 26.9\,\%
goodput tax at 1\,MiB messages, a factor of about 16, because each loss forces
the replay of the requester's outstanding window (about 16 packets on average,
bounded by the 256 packets of the message) rather than of the lost packet
alone. The counter semantics of
\anom{19}: \ctr{packet\_seq\_err} increments once per contiguous loss burst and
not once per lost packet, giving the measured ratio of about 73 between packets
the receiver discarded and sequence errors the sender recorded. And \anom{09},
the firmware-dependent \ctr{local\_ack\_timeout\_err}, selected by
\csr{FW\_VARIANT}. The local acknowledgement timeout must be finite and
configurable: the campaign's most costly instrument defect was a default
encoding that means infinite, under which the first dropped packet stalls a
requester permanently, silently, at zero rate and with no error.

\textbf{Evidence.} Go-back-N and the timeout are \evdoc{}; adaptive
retransmission is \evdrv{}; the selective-repeat window is \evdecl{} and
disabled.

\textbf{Golden model.} The requester transport.

{\raggedright
\textbf{Registers.} \ctr{packet\_seq\_err}, \ctr{roce\_adp\_retrans},
\ctr{local\_ack\_timeout\_err}, \csr{RETRY\_EXCEEDED}, \csr{RTO\_NS};
extension half \csr{UNACKED\_WINDOW\_BYTES}, \csr{REPLAY\_BYTES},
\csr{LOSS\_BURST\_LEN\_HIST}.\par}

\subsubsection{B14: Notification point}

\textbf{Function.} Decide that the endpoint is congested and say so. Two
inputs: a congestion-experienced bit on an arriving packet, which is the
textbook input, and B10's own ingress congestion signal, which is the input
that the measured deployment actually exercises. Generate a congestion
notification packet back to the source queue pair, rate limited to at most one
per configured minimum interval per queue pair.

\textbf{State.} Per queue pair the timestamp of the last notification sent; the
congestion flag; the notification transmit queue.

\textbf{Interfaces.} Congestion indication from B10; the CE flag from B11;
notification packets to B7; counters to B19.

\textbf{Must reproduce.} \anom{16}, the finding that reframes this block. On the
measured fabric the switch never marks: zero congestion-experienced packets in
$6.7 \times 10^{8}$, at two differentiated-services classes, with the switch
buffer full and dropping 6.36\,\%, and with every packet markable. So the
primary input is the endpoint's own ingress congestion, and a fabric that never
marks is the normal case rather than a degenerate one. The rates are the
acceptance data: \ctr{np\_cnp\_sent} at 38 per second at idle rising to 2262
per second exactly during reliable fan-in, which is 283 notifications per
second per congested queue pair, or one per 3.54\mS{}, a factor of 70 under the
one-per-50\uS{} pacing bound the algorithm allows. Below saturation a lone
paced flow must emit exactly zero. And \anom{07}: the marking counter is
present, named as the real device names it, and inert.

\textbf{Evidence.} The congestion-experienced path is \evdoc{}. The
ingress-congestion path is \evcal{}. The notification threshold and the minimum
interval are \evcal{} and fitted over declared candidate grids, because the
vendor register block that would hold them is unreadable on any campaign host.

\textbf{Golden model.} Rate control, notification point.

{\raggedright
\textbf{Registers.} \ctr{np\_cnp\_sent},
\ctr{np\_ecn\_marked\_roce\_packets} (inert mirror),
\csr{CNP\_MIN\_INTERVAL\_NS}; extension half \csr{NP\_CONGESTION\_EVENTS},
\csr{NP\_THRESHOLD\_BYTES}, \csr{NP\_SOURCE} (which input fired).\par}

\subsubsection{B15: Reaction point}

\textbf{Function.} Per-queue-pair rate control. On a notification, update
$\alpha$ toward one, cut the current rate, and remember the pre-cut rate as the
target. In the absence of notifications, recover through fast recovery,
additive increase and hyper increase, gated by both a timer and a byte counter.
Publish the resulting rate to B6.

\textbf{State.} Per queue pair $\alpha$, current rate, target rate, stage, the
timer and the byte counter, and the arrival history the stage machine needs. The
state lives in B4's context store and is cached here.

\textbf{Interfaces.} Notifications from B11; rate to B6; context read and write
to B4.

\textbf{Must reproduce.} \anom{18} in full. A rate cut of at least 30\,\%
arriving after 3, 16 and 39\mS{} across three repeats; fair share reached after
5\mS{}, 1.8\,s and 2.3\,s; recovery to at least 95\,\% of the solo rate in $447
\pm 10$\mS{}; additive increase at about 0.107 falling to 0.093\Gbps{} per
millisecond; and a steady two-flow split of 50.18 and 50.13\Gbps{} on a
99.30\Gbps{} wire. The reaction is 2 to 20 times slower and the recovery about
$4.5\times$ slower than the algorithm's published defaults, so the block
carries fitted parameters rather than the defaults, and says so. Two counter
facts belong here too: \ctr{rp\_cnp\_ignored} is exactly zero, which is what
proves the reaction point is enabled rather than merely present; and senders
handle $2.24\times$ the notifications the receiver reports sending, reproducibly
and without explanation (\anom{19}). The architecture publishes both counters
and does not manufacture the ratio.

\textbf{Evidence.} The algorithm is \evdoc{}. Every parameter value is
\evcal{}, fitted, and marked as such.

\textbf{Golden model.} Rate control, reaction point, in front of the transmit
pacer.

{\raggedright
\textbf{Registers.} \ctr{rp\_cnp\_handled}, \ctr{rp\_cnp\_ignored},
\ctr{roce\_slow\_restart}, \ctr{roce\_slow\_restart\_cnps},
\ctr{roce\_slow\_restart\_trans}, \csr{DCQCN\_AI\_STEP},
\csr{DCQCN\_TIMER\_US}, \csr{DCQCN\_BYTE\_COUNTER}, \csr{DCQCN\_RATE\_FLOOR};
extension half \csr{RP\_ALPHA}, \csr{RP\_CURRENT\_RATE},
\csr{RP\_TARGET\_RATE}, \csr{RP\_STAGE}.\par}

\subsubsection{B16: Completion engine}

\textbf{Function.} Tell the host that work has finished. Write entries into
the completion queue with the correct ownership phase, moderate event-queue
entries and interrupts, and reclaim send-queue slots for unsignalled work when
a later signalled completion retires them.

\textbf{State.} Per completion queue the producer index and owner phase;
moderation counters and timers; event-queue state.

\textbf{Interfaces.} Retirements from B12 and B13; AXI4 write to host memory;
interrupt requests to B1; reclaim notifications to B2.

\textbf{Must reproduce.} That unsignalled work produces no completion entry at
all and that its slots are reclaimed by a later signalled completion, which is
the mechanism behind the difference between transport retirement and
provider-visible slot reclamation. That the ownership bit flips on every wrap so
that a poller sees a coherent entry. And the campaign's negative result: no
completion-queue error was ever observed, and completion service was never on
the critical path at these rates.

\textbf{Evidence.} Entry formats, ownership, the 64 and 128 byte profiles and
moderation are \evdoc{}; the reclaim rule is \evdrv{}.

\textbf{Golden model.} The completion-queue write stage, ordered retirement,
signalled and unsignalled reclaim, and owner wrap.

{\raggedright
\textbf{Registers.} \csr{CQES\_WRITTEN}, \csr{CQE\_ERRORS},
\csr{EQ\_EVENTS}, \csr{MSIX\_INTERRUPTS}, \csr{CQ\_MODERATION\_COUNT},
\csr{CQ\_MODERATION\_US}; extension half \csr{CQ\_OCCUPANCY},
\csr{UNSIGNALLED\_RECLAIMED}.\par}

\subsubsection{B17: Loopback and internal switch}

\textbf{Function.} Carry a packet whose destination is this endpoint back into
the receive path without going out to the wire, and arbitrate between that
traffic and wire traffic for the shared processing budget.

\textbf{State.} The budget token bucket; per-source arbitration state.

\textbf{Interfaces.} Loopback-destined packets from B7; packets to B10; a
service-budget interlock with B10.

\textbf{Must reproduce.} \anom{05}, which is one of the cleanest results in the
campaign. With loopback and wire traffic driven into one adapter together,
total ingress caps near 197\Gbps{}; the wire flow keeps 99 to 99.8\,\% of its
solo rate; the loopback flow absorbs the entire shortfall and falls to about
51\,\%; and no PCIe stall counter and no receive-buffer counter moves. Strict
priority for wire traffic, over a shared budget, is the arbitration rule that
produces exactly that split, and the registered hypothesis that it was PCIe
pressure is refuted by the stall counters. The practical consequence, which the
serving simulator needs, is that colocated ranks exchanging over adapter
loopback starve under wire load.

\textbf{Evidence.} \evcal{} throughout: both the budget and the priority order
are inferred from the split.

\textbf{Golden model.} The internal arbiter, registered against an open task and
not yet landed.

{\raggedright
\textbf{Registers.} \csr{INTERNAL\_BUDGET\_BPS}, \csr{LOOPBACK\_PACKETS},
\csr{LOOPBACK\_BYTES}; extension half \csr{LOOPBACK\_GRANTS},
\csr{WIRE\_GRANTS}, \csr{BUDGET\_OCCUPANCY}.\par}

\subsubsection{B18: Management}

\textbf{Function.} The control plane. Hold the global identifier table,
configure port and maximum-transmission-unit attributes, execute queue-pair
transitions, and run the command interface through which all of that arrives.

\textbf{State.} The global identifier table; port attributes; the command
queue.

\textbf{Interfaces.} AXI4-Lite commands from B1; context writes to B4; table
writes to B5; port configuration to B9.

\textbf{Must reproduce.} The configuration surface as measured. Distinct table
slots for RoCE v1 and RoCE v2 entries, with the v2 entry the one used. A port
maximum transmission unit (4096) that differs from the network-device maximum
transmission unit (9000), because both exist and they are not the same thing.
A priority mapping under which all traffic lands on priority zero, and a
non-identity priority-to-traffic-class map, both of which the campaign found on
every node and neither of which had any effect once priority flow control was
off.

\textbf{Evidence.} \evdoc{}.

\textbf{Golden model.} Not modelled; the C model's caller configures the
endpoint directly. This block is \evdecl{}.

{\raggedright
\textbf{Registers.} \csr{GID\_TABLE\_SIZE}, \csr{ACTIVE\_MTU},
\csr{PORT\_STATE}, \csr{PRIO\_TC\_MAP}, \csr{TRUST\_MODE},
\csr{COMMANDS\_EXECUTED}, \csr{COMMAND\_ERRORS}.\par}

\subsubsection{B19: Control registers, counters and timestamping}

\textbf{Function.} Own the monitoring ring, the versioned map, the free-running
timestamp, the snapshot mechanism and the compatibility behaviour described in
\S\ref{sec:arch:csr}. Every other block's registers are endpoints on this
block's ring.

\textbf{State.} The shadow counter banks; the snapshot timestamp; the map
version and per-block versions.

\textbf{Interfaces.} AXI4-Lite ring to all blocks; the broadcast timestamp and
snapshot strobe; the AXI4-Lite slave that B1 drives.

\textbf{Must reproduce.} Counter \emph{semantics} and not merely counter
names. There are three of them. \anom{07}: the marking counter
\ctr{np\_ecn\_marked\_roce\_packets} is present and permanently zero while
congestion control runs. \anom{09}, the firmware-selected behaviour of
\ctr{local\_ack\_timeout\_err}. \anom{19}, the whole set: sequence errors
counting bursts, ingress discards counting packets exactly, switch loss
appearing in neither, the receive-buffer counter never moving, and the
notification ledger that does not balance. These are facade behaviours with no
datapath effect, and a tool written against the real device must behave
identically against this one, which is the entire point of the mirror half of
the map.

\textbf{Evidence.} Counter names, widths and reset behaviour are \evdoc{} from
the kernel's counter documentation; the semantics that contradict the names are
measured.

\textbf{Golden model.} The observable-state facade and the firmware counter
variant field.

{\raggedright
\textbf{Registers.} \csr{MAP\_VERSION}, \csr{SNAPSHOT},
\csr{TIMESTAMP\_LO}, \csr{TIMESTAMP\_HI}, \csr{WRAP32\_COMPAT},
\csr{FW\_VARIANT}, \csr{RING\_STOPS}, \csr{RING\_ERRORS}.
\par}

\subsection{Expandability}
\label{sec:arch:expand}

Table~\ref{tab:params} lists the parameters that a different target changes.
None of them is buried in a datapath width; each is a top-level generic that
propagates into buffer depths, counter widths and the register map's
\csr{CAPABILITIES} words.

\begin{table}[t]
  \centering
  \small
  \caption{Parameters. The U55C column is the first build; the range is what
    the architecture is designed to accept without restructuring.}
  \label{tab:params}
  \footnotesize
  \setlength{\tabcolsep}{3.5pt}
  \begin{tabular}{@{}ll>{\raggedright\arraybackslash}p{0.42\columnwidth}@{}}
    \toprule
    Parameter & U55C build & Range \\
    \midrule
    Link rate per port  & 100\Gbps{}    & 100 to 400\Gbps{} \\
    Ports               & 2             & 1 to 4 \\
    Datapath width      & 512\,b        & 512 to 2048\,b \\
    Path MTU            & 4096\,B       & 256 to 4096\,B \\
    Queue pairs         & 16\,k on chip & 1\,k to 1\,M with backing \\
    Memory regions      & 64\,k         & 1\,k to 1\,M with backing \\
    Ingress depth       & \evcal{} fit   & any \\
    Internal budget     & 197\Gbps{}    & scales with the datapath \\
    Context backing     & HBM2          & HBM, DDR or on-die SRAM \\
    Retransmission      & go-back-N     & plus optional selective repeat \\
    \bottomrule
  \end{tabular}
\end{table}

\textbf{Rate.} Moving from 100 to 400\Gbps{} widens the datapath rather than
raising the clock: 2048 bits at 322.266\,MHz carries 660\Gbps{}. The blocks
that must change are the packetizer, which segments per beat rather than per
byte, and the ingress meter, whose service rate is a parameter. The
congestion-control constants do \emph{not} scale linearly, and the model marks
every ConnectX-7 field derived by scaling as \evdecl{} rather than pretending it
was measured.

\textbf{Scale.} Queue-pair and memory-region counts above the on-chip cache
capacity are served from high-bandwidth memory, and above that from host memory
over the same translation path the real device uses. The architecture's
position here is that a cache miss must be attributable: a lookup that misses
publishes the miss in its own counter and the service time it cost, so a scale
study can tell a context miss from a translation miss, which the measured
device cannot.

\textbf{Application-specific path.} Three substitutions and nothing else:
high-bandwidth memory becomes on-die static RAM with the same AXI4 interface;
the Ethernet and PCIe soft wrappers become hard macros; and the monitoring ring
keeps its width but loses the FPGA-only extension half if area demands it, with
\csr{CAPABILITIES} advertising the loss. The register-transfer-level sources
and the register map are unchanged, which is the whole point of putting the
observation path on its own interconnect.
